\chapter{绪论}

本章从起源背景与发展现状等维度，对AI任务调度系统进行了阐述，并深入剖析了现有技术方案存在的问题。针对这些问题，本章进一步调研了国内外及开源领域技术方案在发展现状中存在的缺陷。为解决这些缺陷，本系统围绕AI相关特性适配、AI任务可执行性判断以及AI多任务资源互斥检测等方面，提出了相应解决方案，以使本系统能够契合AI发展的需求。 

\section{选题背景及意义}

\subsection{AI大模型背景介绍}

习近平总书记指出，“人工智能是新一轮科技革命和产业变革的重要驱动力量”\cite{lu2025powermodel}。AI大模型是在数据爆炸式增长，传统机器学习模型处理大规模数据效率低、应对复杂任务能力不足，以及硬件计算能力提升的场景下被提出。其提出解决了诸多问题，在自然语言处理上提升语言理解与生成能力；在计算机视觉领域可处理复杂场景分析；具备知识融合与推理能力以辅助决策；还具有通用性和迁移性，能提高开发效率。 

AI从神经网络发展到大模型，是一个持续演进与突破的历程\cite{subakan2021attention, radford2018improving, devlin2019bert, radford2019language, chowdhery2023palm, rae2021scaling, bubeck2023sparks}。2017年，Google提出的Transformer架构\cite{subakan2021attention}成为关键转折点，其通过自注意力机制和并行计算特性，为现代大型语言模型奠定基础；2018年OpenAI推出的GPT-1验证了“预训练-微调”范式\cite{radford2018improving}，谷歌BERT推动自然语言理解进步\cite{devlin2019bert}，后续GPT-2\cite{radford2019language}、GPT-3进一步展现生成与泛化能力；2020年后大模型参数规模指数级增长，2022年谷歌PaLM\cite{chowdhery2023palm}、DeepMind Gopher\cite{rae2021scaling}等问世，2023年GPT-4\cite{bubeck2023sparks}实现多模态，更多厂商加速研发，多模态模型不断涌现，AI应用场景愈发广泛。

人工智能大模型凭借强大的学习与处理能力，已深度融入工程、网络舆论、医疗等众多领域\cite{geng2025application, he2025framework, zhang2022medical}。耿潇\cite{geng2025application}等将大模型应用在电力系统造价预测、工程审核中。何巍\cite{he2025framework}等将大模型应用于网络谣言治理领域。张华丽\cite{zhang2022medical}等将大模型应用于医疗问答领域。广泛的应用不仅重塑了行业生态，也为社会经济发展注入强劲动力。随着人工智能技术的持续演进，用户对其需求呈现爆发式增长，这促使各技术公司竞相投入资源，加速人工智能大模型的开发与测试进程。

\subsection{AI多任务调度的研究意义}

在AI产品的生命周期中，存在着许多任务需要调度执行。例如开发与测试环节，程序员需要编写任务进行初步的自测和反馈；AI模型上线后，系统也有启用、禁用模型接口等任务需要调度。通常，程序员会将这些工作编写成任务脚本，然后启动脚本来执行。然而随着AI产业蓬勃发展，人工调度任务早已不能满足行业需要，尤其是在大量的开发与应用场景中，AI任务数量庞大，要对每个任务进行控制、发起并调度，更是一个棘手的难题。 

为解决上述难题，现阶段工业界主要采用两种方法：一是人工编排发起任务，二是通过脚本控制任务队列来分配任务。然而，这两种方法存在诸多问题。在非结构化参数的手动标注方面，由于缺乏统一标准，标注人员理解各异，极易出现标签不一致、标注错误等情况，进而影响AI任务的执行效果。手工安排AI模型测试任务时，若未能充分考量任务间的依赖关系，可能致使任务执行顺序错乱，不仅会使测试结果失效，甚至可能引发系统冲突。此外，还容易因疏忽导致重要任务延迟执行。从模型调用层面看，人工频繁输入指令调用AI大模型，可能因指令格式错误、参数设置不当，致使模型无法响应或输出无效结果。在批量调用时，人工难以精准把控调用频率，可能触发模型使用限制，进而影响工作进度。而且，人工操作缺乏自动化监控机制，一旦出现异常情况，无法及时察觉与处理，问题不断累积将进一步扩大损失。在传统模式下，每个操作环节都需程序员手动执行，操作步骤既繁琐又机械。大量时间耗费在等待任务完成、重复操作以及错误排查上，难以实现多任务并行处理，无法充分利用计算资源，使得整体效率远低于自动化调度系统，严重制约了人工智能产品的快速迭代。
AI的快速发展与传统调度手段效率低下之间的矛盾，构成本工作的核心研究背景，这一矛盾可以拆解为三个亟待解决的子问题：一，当前系统与AI的适配性欠佳，对AI相关任务的数据处理能力不足，无法有效完成数据纠错与类型分析；二，无法依据内外部条件自动决策任务的执行与否，过度依赖人工判断，导致系统效率大幅降低；三，缺乏对任务间资源互斥情况的智能分析能力，导致任务往往只能串行执行，造成计算资源的严重浪费。

AI任务对调度系统的适配性要求较高，同时需具备感知内外部条件以判断任务可执行性的能力，并且要充分考虑任务间资源互斥的情况。由此可见，现阶段的AI多任务调度方案已无法满足AI的发展需求。

\section{国内外发展现状}

\subsection{任务调度技术及其对AI的适配现状}

自动化任务调度与监控系统的发展可追溯到20世纪60年代英国曼彻斯特大学Atlas计算机的底层资源管理与任务调度机制\cite{kilburn1961atlas}。当时随着计算机技术兴起，高校和企业尝试将自动化技术应用于生产和管理，该技术应运而生。早期系统以批处理作业调度为主，核心目标是解决计算机系统中简单、单一的任务分配问题，其架构设计围绕 “固定流程执行” 展开，无法适配AI大模型高算力、长周期、多模态数据交互的特性。

20世纪80年代，随着个人计算机普及和互联网发展，该技术逐渐从单一系统内部调度扩展到跨平台、跨网络的分布式调度\cite{birrell1984implementing}，虽开始融入人工智能、大数据等技术，但此时的AI应用多为基础算法（如规则引擎、简单机器学习模型），传统调度系统的优化方向仍聚焦于“提升多任务并发效率”，未就AI大模型的核心特点进行针对性设计。例如，系统资源分配采用“平均分配”或“按固定阈值分配”模式，无法满足AI大模型训练时对GPU、显存等专用资源的突发性、高需求；数据处理接口仅支持结构化数据，难以兼容AI大模型所需的文本、图像、音频等多模态数据，进一步凸显“未针对AI大模型任务进行优化”的局限。

21世纪以来，随着云计算、物联网、大数据等新兴技术发展，自动化任务调度与监控系统虽向智能化、可视化、移动化方向发展，且技术成熟度显著提升，但传统任务调度系统并未针对AI大模型任务进行优化的问题愈发凸显。在实际应用中，这一局限导致两大核心矛盾：一是资源利用率低下，AI大模型训练需持续占用大量算力资源，而传统系统缺乏“资源动态伸缩”机制，闲置时造成算力浪费，峰值时又因资源不足导致任务中断；二是调度效率不足，AI大模型任务常需跨节点、跨集群协同，传统系统的任务依赖管理机制无法应对其复杂的依赖关系，易出现任务阻塞，例如在制造业智能质检场景中，AI大模型推理任务因传统调度系统无法优先分配算力，导致检测延迟增加，影响生产节奏。

开源的任务调度框架，如Airflow\cite{finnigan2021}、Oozie\cite{bandi2018}等，具有较高的灵活性和可扩展性，吸引了众多开发者参与贡献和二次开发。然而，由于开源项目的开发者众多，代码风格和质量参差不齐，缺乏统一的维护和管理，导致其稳定性和可靠性存在一定风险。同时，开源项目的功能往往较为基础，对于AI多任务场景也没有特别的适配。此外，开源AI任务调度系统通常具有较高的技术复杂性，对于初学者和小型团队来说，学习和掌握这些系统的使用方法需要花费大量的时间和精力。例如，Airflow的工作流定义较为复杂，需要开发者对其有深入的理解才能正确地进行任务调度配置；Oozie有着较好的性能，但配置复杂\cite{mei2022}。

AI作为新兴技术，其多任务调度问题尚未引起传统调度系统、算法及框架的关注，现有相关工作难以契合AI的发展需求。本系统的任务调度层正是为解决这一问题而设计。本系统应降低AI多任务调度的操作难度，简便调度过程；智能判断任务可执行性，节约人力成本；识别多任务场景下的资源互斥情况，提高任务并发度，进而提升整体运行效率。

\subsection{任务可用性检测机制研究现状}

工程上，任务可用性检测主要通过逻辑引擎来实现。早期逻辑引擎以经典单调推理为核心\cite{liao2008defeasible}，依赖预定义规则与静态知识完成基础检测，仅适用于结构化场景的简单任务，在动态环境中易出现检测偏差，并不适用于本文所要面对的内外部条件变化检测场景。随后非单调推理技术兴起，麦卡锡限制\cite{mccarthy1980circumscription}、赖特默认理论\cite{reiter1980logic}等框架与现代模型论融合，使逻辑引擎具备动态调整检测逻辑、处理可废止结论的能力，提升了复杂场景下的判断准确性，但是应用在本场景下显得过于笨重，不利于系统的维护。再到演绎算法向动态自适应演进，通过优化推理策略提升了高复杂度任务的检测效率，不过受限于静态条件依赖等问题，其在AI多任务系统中仍有不足。

现有逻辑引擎虽在特定领域推理与优化中展现价值，但存在显著功能局限：其推理过程依赖预定义规则与静态知识框架，缺乏对AI多任务系统中任务目标、资源约束、环境动态性的综合感知与建模能力，无法有效分析任务执行所需的前提条件、资源匹配度及流程兼容性，因此无法作为AI多任务系统判断任务可执行性的核心支撑。

\subsection{资源互斥检测机制研究现状}

早期，经典软件互斥算法，如Dekker算法与Peterson算法问世，旨在借助共享变量及既定逻辑规则，实现多进程对临界资源的互斥访问。然而，这些算法存在局限性。以Peterson算法为例，其仅适用于双进程场景，扩展至多进程时逻辑复杂，且软件方法普遍存在 “忙等” 特性，易造成CPU资源浪费。在AI场景中，资源需求动态波动且多进程并发访问频繁，此类算法扩展性不足与资源浪费问题被放大，难以适配AI资源的灵活调度需求。

硬件方法\cite{bershad1992fast}利用处理器提供的原子指令，如“测试并设置”“交换”等，实现临界区访问控制，具备通用性、高效性与安全性。但同样存在 “忙等” 问题，且其设计主要针对传统同构计算资源，面对AI场景中GPU、TPU等异构资源的混合调度，缺乏针对性优化，无法充分匹配AI资源特殊性。虽实际系统常结合硬件原子指令与“阻塞、唤醒”机制解决“忙等”问题，却仍未突破传统资源调度框架，难以满足AI资源在任务执行中动态调整资源分配的需求。

在分布式系统领域，Lamport\cite{lamport2019time}提出基于时间戳的互斥算法，通过逻辑时钟为进程请求分配时间戳，依此决定临界资源访问权限，解决了分布式系统进程同步难题。G.Ricart\cite{ricart1981optimal}算法等也在分布式资源分配中得以研究与应用。然而，这些分布式互斥算法多基于固定的资源请求优先级与访问逻辑，在AI分布式训练中，资源请求与任务进度、模型复杂度动态相关，不同训练阶段对算力、内存需求差异大，传统算法难以据此动态调整资源互斥策略，对AI资源特殊性适应性不足。

近年来，线程组规模扩大、互斥资源复杂，线程组死锁问题凸显。国内学者针对分布式系统特点与实际需求，提出诸多改进的死锁检测、预防算法及适应性与性能更佳的互斥算法，循环编码的互斥请求集产生算法\cite{DZXU200508011}、锁链长度（LCL）的分布式互斥检测算法\cite{yang2023lcl}等，提高了系统并发性能与资源利用率。

由于现有算法未能全面考虑AI资源的独特属性，这些属性表现为资源需求的动态变化，同时面临资源碎片化以及多任务间依赖关系复杂等难题。在此情况下，现有算法在协调不同类型AI任务的资源互斥与效率优化之间的平衡上存在不足，这充分表明传统资源互斥算法难以适应AI资源的特殊性。因此，本系统有必要设计一套专门适用于AI任务，特别是与GPU等AI特有资源可用性相关的互斥检测算法。 

\section{本文主要工作}

针对当前主流调度系统在AI大模型多任务场景下存在的适配性差、自动化程度低及计算资源浪费等问题，本文设计并实现了一套面向AI多任务场景的资源调度系统。系统通过API接入层、任务调度层与任务执行层的分层架构设计，系统性地完成了以下工作：

\begin{enumerate}
    \item 针对AI适配性与数据处理能力不足的问题，在API接入层与任务执行层建立了大模型任务参数画像与异构数据映射机制。通过设计数据类型过滤层与自动化纠正逻辑，实现了在复杂参数条件下AI任务数据的自动纠错与类型推断，显著降低了系统的操作复杂度与数据处理门槛。
    \item 针对任务决策过度依赖人工判断的问题，在任务调度层构建了基于逻辑引擎的任务可用性检测机制。系统通过实时感知内外部环境（如系统负载、运营策略等），利用配置化规则引擎实现了调度指令的自动下发与执行决策，消除了任务启动过程中的人工干预。
    \item 针对计算资源互斥导致的任务串行效率低下问题，设计了针对AI特有资源的互斥智能分析算法。通过建立细粒度的“临界资源表”与权重轮询负载均衡策略，实现了对GPU等计算资源的精准识别与动态分配，推动系统由传统串行执行向高效并发调度转变。
    \item 完成了系统的整体集成与多维度性能验证。本文通过功能性测试验证了分层架构在复杂业务逻辑下的正确性，并利用Hey等工具在生产环境（腾讯TEG业务场景）完成了性能与稳定性评估。实验结果表明，系统每月可稳定支撑约200个AI任务运行，显著提升了资源利用率并大幅降低了运维人力投入。
\end{enumerate}

\section{本论文的组织结构}

本论文共分为七章，各章节的具体内容与研究逻辑安排如下：

第一章，绪论。主要阐述人工智能大模型背景下任务调度系统的研究意义，深入分析现有编排模式在AI适配性、决策自动化及资源冲突处理方面的不足，引出本文拟解决的三个核心子问题，并概述本文的研究内容与创新点。

第二章，相关技术简介。介绍实现本系统所需的Python、Celery分布式任务队列、FastAPI框架、tRPC协议及PostgreSQL数据库等技术栈。重点阐述选用这些技术如何支撑AI多任务场景下灵活的数据处理与高效的任务流转。

第三章，需求分析。针对AI多任务场景下的三大核心矛盾展开深度需求剖析。分别明确API接入层在异构数据纠错、任务调度层在自动化可用性检测以及执行层在资源互斥识别等方面的具体功能需求与性能目标。

第四章，概要设计。确立系统的三层架构方案，包括API功能接入层、任务调度层与任务执行层。通过层级化设计建立系统基本原则，初步构建解决AI特性适配、任务自动调度及资源分类隔离的架构模型。

第五章，详细设计与实现。本章详细描述了各模块的具体实现。API接入层重点实现权限鉴定与数据校验，解决AI异构数据的适配难题；任务调度层通过状态机、逻辑引擎及资源互斥算法，实现调度自动化与高效并发；任务执行层明确参数传递、任务编辑与异常捕获机制，保障执行鲁棒性。

第六章，系统测试与分析。通过功能性测试验证系统对三大子问题的解决情况，确保API纠错、自动调度与资源互斥处理符合预期。同时开展非功能性测试，评估系统在生产环境下的性能表现、稳定性和兼容性。

第七章，结论与展望。总结本文在攻克AI多任务调度难题方面的研究贡献，阐述系统在提升调度效率、节省人力成本方面的实际应用价值，并针对未来适配多样化临界资源及完善AI智能报表等方向做出展望。

% \makeatletter
% \@openrightfalse
% \makeatother
